% PROPOSTA (P-09) — versao <=500 palavras, norma SiBi/UFPR.
% NAO substitui 0-iniciais/resumo.tex. O corte e escolha editorial do autor.
%
% O QUE ESTA VERSAO E, medido e nao alegado: CONDENSACAO FIEL, nao recorte.
% Das 19 frases, 3 sao verbatim do resumo completo e 16 foram reescritas —
% frases fundidas e oracoes comprimidas para caber no teto. Uma versao
% anterior deste cabecalho dizia "removido, nunca reescrito"; era FALSO, e a
% medicao do proprio autor desta proposta o refutou.
%
% O que NAO mudou, e isso continua provado por script: nenhum numero novo
% (o conjunto de numeros do curto e SUBCONJUNTO do completo), nenhuma
% alegacao que nao esteja na versao completa, e o corte e identico em PT e EN.
\begin{resumo}

A classificação de textos curtos em português depende de grandes volumes de rótulos, e essa rotulagem é o principal custo desses sistemas. O caso estudado é o das descrições de produtos de notas fiscais eletrônicas, com vocabulário abreviado e centenas de categorias desbalanceadas. Esta tese propõe e avalia o FALCO (\textit{Framework} de Aprendizado Ativo com LLM para texto CurtO), que ataca esse custo em três frentes: composição informada do conjunto inicial, seleção ativa por incerteza e substituição parcial do anotador humano por oráculos baseados em modelos de linguagem de grande porte (LLMs), em fases de custo crescente. A investigação usa um conjunto público auditado de 250.365 descrições em 621 categorias. O protocolo é pré-registrado, com partições imutáveis, análise pareada (McNemar, Wilcoxon) e artefato rastreável para cada número. Cinco resultados sustentam a tese. Primeiro, a composição do conjunto inicial importa mais quanto menor o orçamento, e a heurística proposta, o DRI-SL, supera sem usar nenhum rótulo o melhor indivíduo de um algoritmo genético supervisionado em todos os tamanhos de 100 a 5.000. Segundo, oráculos LLM \textit{zero-shot} atingem $77$--$83\%$ de acurácia no espaço fechado de 621 categorias, a custos de US\$~0,035 a US\$~0,92 por mil rótulos, com Macro F1 superior ao de um classificador supervisionado leve treinado com a base completa. Terceiro, o instrumento de medição é parte do resultado: sem restrição de saída, $6{,}8\%$ das respostas viram falsos erros, e o mesmo modelo gratuito diverge significativamente de si mesmo entre dois provedores de \textit{serving} ($p<0{,}001$). Quarto, as estratégias de incerteza superam a seleção aleatória na significância máxima que 8 sementes permitem ($p=0{,}0078$), e a vantagem sobrevive a ruído uniforme de oráculo até $\varepsilon=0{,}4$, faixa que cobre o erro dos LLMs reais. Nenhum oráculo alcançou o critério pré-registrado de $85\%$ de acurácia, o que tornou o aprendizado com rótulos ruidosos o cenário central. Quinto, o framework foi executado ponta a ponta com oráculo LLM gratuito, a custo monetário zero, e validado com o BERTimbau em três sementes. A hipótese central era obter pelo menos $95\%$ da acurácia da supervisão completa do \textit{pool} de referência com no máximo 34.724 rótulos. O veredito é duplo. \textbf{O critério é atingível dentro do teto, e é cruzado com 20 mil rótulos ($8{,}6\%$ da população) nas três sementes, em varredura \textit{post hoc} com rótulos de gabarito. A configuração executada, porém, parou por estagnação aquém do piso, com 11.936 rótulos ($5{,}2\%$)}. A decomposição pareada inocenta o ruído do oráculo e a seleção, e atribui a perda ao critério de parada herdado do classificador leve. As limitações são declaradas: os resultados provêm de uma única base, do próprio domínio de estudo, e a análise de custo apoia-se em razões entre oráculos, não em preços absolutos. A tese entrega ainda a métrica LCE (\textit{Learning Curve Efficiency}), a biblioteca aberta \texttt{activelearning} e a ferramenta FlowBuilder, que tornam o protocolo reutilizável.

\end{resumo}
